% ==============================================================================
% lab2_forensics_disque.tex — UE SIS589
% Travaux Pratiques 2 : Acquisition forensique, Autopsy, Timeline, Rapport
% ==============================================================================

\chapter{TP~2~— Acquisition Forensique, Analyse de Disque et Rapport}
\label{lab:forensics}

\epigraph{\itshape « Chaque fichier raconte une histoire.
  L'investigateur lit entre les secteurs. »}
         {--- SANS FOR500}

% ── Présentation ──────────────────────────────────────────────────────────────
\section*{Présentation du TP}

\begin{table}[H]
\centering\small
\rowcolors{2}{TableauPair}{TableauImpair}
\begin{tabular}{p{3.5cm}p{10cm}}
\toprule
\textbf{Durée} & 6~heures (2~h acquisition + 4~h analyse + rédaction) \\
\textbf{Objectifs} & Acquérir une image disque bit à bit~; maintenir la
  chaîne de custody~; analyser les artefacts Windows avec Autopsy et EZTools~;
  analyser un dump mémoire avec Volatility~3~; construire une timeline
  forensique et rédiger un rapport d'investigation \\
\textbf{Pré-requis} & Chapitres~5, 6 et~7~; connaissance de base de Linux \\
\textbf{Infrastructure} & Station SIFT Workstation fournie (VM)~;
  image disque \texttt{srv-web-01.dd} et dump mémoire
  \texttt{srv-web-01.lime} disponibles sur le partage TP~;
  bloqueur en écriture USB disponible en salle \\
\textbf{Livrable} & Rapport d'investigation forensique complet
  (20~pages max, format académique) \\
\bottomrule
\end{tabular}
\end{table}

\textbf{Contexte du cas.}
Le 12~avril~2026 à 03h14~UTC, le SIEM de LiZbiZ-SIS a détecté une
intrusion sur \texttt{srv-web-01} (Ubuntu~22.04, serveur Apache~2.4,
IP~: 10.0.1.10). L'alerte initiale signalait un brute force SSH réussi
depuis 203.0.113.42. Vous avez été mandaté pour conduire l'investigation
forensique de ce serveur~; les images vous ont été remises à 05h30~UTC
avec le PV de collecte de l'analyste de permanence.


% ==============================================================================
\section{Partie~A — Vérification des pièces et prise en charge}
\label{lab2:partieA}
% ==============================================================================

\subsection{A1 — Vérification d'intégrité des images}

\begin{lstlisting}[style=StyleTerminal,
  caption={Vérification d'intégrité avant toute analyse},
  label={lst:lab2-hash-verify}]
# ─── Sur la SIFT Workstation — Commencer TOUJOURS par cette étape ────────────
# Copier les images sur la station d'analyse (NE PAS ANALYSER sur le partage)
mkdir -p ~/cases/INC-2026-0412-001/evidence/
cp /mnt/tp-share/srv-web-01.dd       ~/cases/INC-2026-0412-001/evidence/
cp /mnt/tp-share/srv-web-01.lime     ~/cases/INC-2026-0412-001/evidence/
cp /mnt/tp-share/pv_collecte.txt     ~/cases/INC-2026-0412-001/evidence/

# Vérifier les hashes contre ceux du PV de collecte
sha256sum ~/cases/INC-2026-0412-001/evidence/srv-web-01.dd
sha256sum ~/cases/INC-2026-0412-001/evidence/srv-web-01.lime

# Comparer manuellement avec le PV de collecte
# Les deux valeurs DOIVENT être identiques

# Si mismatch : ARRÊT IMMÉDIAT — contacter l'enseignant
# et documenter la divergence dans votre rapport
\end{lstlisting}

\begin{boxexo}
\textbf{Question A1.1~:} Lisez le fichier \texttt{pv\_collecte.txt}.
Identifiez les 5~éléments obligatoires d'un PV de collecte conforme
à ISO~27037. Sont-ils tous présents~? Signaler toute lacune.

\medskip\textbf{Question A1.2~:} Les hashes SHA-256 calculés
correspondent-ils aux valeurs du PV~? Documentez le résultat dans
votre rapport (captures d'écran des deux commandes).
\end{boxexo}

\subsection{A2 — Montage de l'image en lecture seule}

\begin{lstlisting}[style=StyleTerminal,
  caption={Montage de l'image disque en lecture seule sur SIFT},
  label={lst:lab2-montage}]
# ─── Identifier les partitions dans l'image ───────────────────────────────────
cd ~/cases/INC-2026-0412-001/evidence/
mmls srv-web-01.dd
# Résultat attendu :
# 000: Meta    0000000000 0000000001 0000000001 Primary Table (#0)
# 001: -----   0000000000 0000002047 0000002048 Unallocated
# 002: 0x83    0000002048 0002099199 0002097152 Linux (ext4) <-- PARTITION
# 003: 0x8E    0002099200 0004194303 0002095104 Linux LVM

# Monter la partition principale en lecture seule
# offset = secteur début × 512 = 2048 × 512 = 1048576
sudo mkdir -p /mnt/forensics/srv-web-01/
sudo mount -o ro,loop,offset=1048576 \
  ~/cases/INC-2026-0412-001/evidence/srv-web-01.dd \
  /mnt/forensics/srv-web-01/

# Vérifier le montage
ls /mnt/forensics/srv-web-01/
mount | grep forensics  # Confirmer "ro" (read-only)

# Vérifier que le hash du disque n'a pas changé
sha256sum ~/cases/INC-2026-0412-001/evidence/srv-web-01.dd
# Doit être IDENTIQUE à la valeur initiale
\end{lstlisting}

\begin{boxexo}
\textbf{Question A2.1~:} Pourquoi utilise-t-on l'option \texttt{ro}
(read-only) au montage~? Quels risques encourrait-on en montant
l'image en lecture-écriture~?

\medskip\textbf{Question A2.2~:} La commande \texttt{mmls} révèle
la structure de la table de partitions. À quoi correspond la zone
\texttt{Unallocated} entre le secteur~0 et~2047~?
\end{boxexo}


% ==============================================================================
\section{Partie~B — Analyse des artefacts Linux}
\label{lab2:partieB}
% ==============================================================================

\subsection{B1 — Analyse des journaux d'authentification}

\begin{lstlisting}[style=StyleTerminal,
  caption={Analyse d'auth.log sur l'image montée},
  label={lst:lab2-authlog}]
# ─── Journaux d'authentification ─────────────────────────────────────────────
# Tous les événements SSH de la nuit du 12 avril
grep "Apr 12" /mnt/forensics/srv-web-01/var/log/auth.log | \
  grep -E "(Failed|Accepted|Invalid)" | \
  awk '{print $1,$2,$3,$9,$11}' | sort

# Top IP sources d'échecs SSH
grep "Failed password" /mnt/forensics/srv-web-01/var/log/auth.log | \
  awk '{print $(NF-3)}' | sort | uniq -c | sort -rn | head -10

# Identifier la première authentification réussie
grep "Accepted" /mnt/forensics/srv-web-01/var/log/auth.log | head -5

# ─── Historique des connexions (wtmp) ─────────────────────────────────────────
last -f /mnt/forensics/srv-web-01/var/log/wtmp | head -20

# ─── Historique bash de chaque utilisateur ────────────────────────────────────
for user in $(ls /mnt/forensics/srv-web-01/home/); do
  histfile="/mnt/forensics/srv-web-01/home/$user/.bash_history"
  if [ -f "$histfile" ]; then
    echo "=== $user ==="
    cat "$histfile"
  fi
done
cat /mnt/forensics/srv-web-01/root/.bash_history 2>/dev/null
\end{lstlisting}

\begin{boxexo}
\textbf{Travail~B1~:} Remplissez un tableau chronologique des
événements SSH~: timestamp, IP source, utilisateur, résultat
(succès/échec). Répondez~:
(a)~Combien d'échecs avant le premier succès~?
(b)~À quelle heure exacte le succès a-t-il eu lieu~?
(c)~Quel compte a été compromis~?
(d)~Quelles commandes suspectes apparaissent dans les historiques bash~?
\end{boxexo}

\subsection{B2 — Recherche de fichiers suspects et persistances}

\begin{lstlisting}[style=StyleTerminal,
  caption={Recherche de fichiers malveillants et persistances},
  label={lst:lab2-suspects}]
# ─── Fichiers créés/modifiés entre 03h00 et 05h00 le 12/04/2026 ──────────────
find /mnt/forensics/srv-web-01/ \
  -newer /mnt/forensics/srv-web-01/etc/passwd \
  -not -newer /mnt/forensics/srv-web-01/etc/shadow \
  -type f 2>/dev/null | \
  grep -v "/proc\|/sys\|/run" | sort

# Alternative avec touch et timestamps
# Créer des fichiers repères avec les timestamps de la fenêtre d'incident
# puis utiliser find -newer

# ─── Fichiers dans les répertoires temporaires ────────────────────────────────
find /mnt/forensics/srv-web-01/tmp \
     /mnt/forensics/srv-web-01/var/tmp \
     /mnt/forensics/srv-web-01/dev/shm \
  -type f -ls 2>/dev/null

# ─── Webshells potentiels dans le DocumentRoot Apache ─────────────────────────
find /mnt/forensics/srv-web-01/var/www/ -name "*.php" -ls 2>/dev/null
grep -rl "eval\|base64_decode\|system\|exec\|passthru\|shell_exec" \
  /mnt/forensics/srv-web-01/var/www/ 2>/dev/null

# ─── Mécanismes de persistance ────────────────────────────────────────────────
# Crontabs
for user in root $(ls /mnt/forensics/srv-web-01/home/); do
  cfile="/mnt/forensics/srv-web-01/var/spool/cron/crontabs/$user"
  [ -f "$cfile" ] && echo "=== Cron $user ===" && cat "$cfile"
done

# Authorized_keys (backdoors SSH)
find /mnt/forensics/srv-web-01/ -name "authorized_keys" 2>/dev/null | \
  xargs -I{} sh -c 'echo "=== {} ==="; cat {}'

# Services systemd non standards
find /mnt/forensics/srv-web-01/etc/systemd/system/ \
  -name "*.service" 2>/dev/null | sort

# LD_PRELOAD (rootkit)
cat /mnt/forensics/srv-web-01/etc/ld.so.preload 2>/dev/null
\end{lstlisting}

\begin{boxexo}
\textbf{Travail~B2~:} Documentez chaque artefact suspect trouvé~:
chemin complet, timestamps MACE (stat), contenu (si lisible),
TTP MITRE correspondant. Utilisez un tableau à 4~colonnes
(chemin, timestamps, contenu, TTP).
\end{boxexo}

\subsection{B3 — Récupération de fichiers supprimés}

\begin{lstlisting}[style=StyleTerminal,
  caption={Carving et récupération de fichiers supprimés},
  label={lst:lab2-carving}]
# ─── The Sleuth Kit : fichiers supprimés ──────────────────────────────────────
mkdir -p ~/cases/INC-2026-0412-001/recovered/
cd ~/cases/INC-2026-0412-001/

# Lister les entrées de répertoire supprimées
fls -r -d ~/cases/INC-2026-0412-001/evidence/srv-web-01.dd | \
  head -50 > recovered/deleted_files_list.txt
cat recovered/deleted_files_list.txt

# Récupérer un fichier spécifique par inode
# (remplacer INODE par le numéro trouvé dans la liste)
# icat evidence/srv-web-01.dd INODE > recovered/fichier_recupere.bin
# sha256sum recovered/fichier_recupere.bin

# ─── Scalpel : carving de l'espace non alloué ────────────────────────────────
mkdir -p recovered/scalpel_output/
scalpel -c /etc/scalpel/scalpel.conf \
        -o recovered/scalpel_output/ \
        evidence/srv-web-01.dd
# Types recherchés : scripts shell, archives, PDF, scripts Python

# ─── bulk_extractor : extraction de patterns ─────────────────────────────────
mkdir -p recovered/bulk_output/
bulk_extractor -R \
  -o recovered/bulk_output/ \
  -e net -e email -e url \
  evidence/srv-web-01.dd
# Analyser : url.txt, email.txt, domain.txt
head -30 recovered/bulk_output/url.txt | grep -v "microsoft\|ubuntu\|debian"
\end{lstlisting}

\begin{boxexo}
\textbf{Travail~B3~:} Identifiez au moins~2~fichiers supprimés
pertinents pour l'investigation. Pour chacun~:
(a)~numéro d'inode~;
(b)~contenu récupéré (si lisible)~;
(c)~hash SHA-256 de la version récupérée~;
(d)~interprétation forensique (à quoi servait ce fichier~?).
\end{boxexo}


% ==============================================================================
\section{Partie~C — Analyse du dump mémoire}
\label{lab2:partieC}
% ==============================================================================

\subsection{C1 — Analyse avec Volatility~3}

\begin{lstlisting}[style=StyleTerminal,
  caption={Analyse du dump mémoire srv-web-01 avec Volatility~3},
  label={lst:lab2-volatility}]
# Toutes les commandes depuis ~/cases/INC-2026-0412-001/
DUMP="evidence/srv-web-01.lime"
V3="python3 /opt/volatility3/vol.py"

# ─── Étape 1 : Identification du système ──────────────────────────────────────
$V3 -f $DUMP linux.info

# ─── Étape 2 : Liste des processus ────────────────────────────────────────────
$V3 -f $DUMP linux.pslist   > analysis/pslist.txt
$V3 -f $DUMP linux.pstree   > analysis/pstree.txt
cat analysis/pstree.txt | head -40

# ─── Étape 3 : Connexions réseau actives au moment du dump ────────────────────
$V3 -f $DUMP linux.netstat  > analysis/netstat.txt
cat analysis/netstat.txt

# ─── Étape 4 : Recherche de processus suspects ───────────────────────────────
# Processus démarrés depuis /tmp, /dev/shm ou des chemins inhabituels
grep -E "(/tmp/|/dev/shm/|/var/tmp/)" analysis/pslist.txt

# Processus sans chemin connu (fileless)
grep -E "^[0-9].*\s-\s*$" analysis/pslist.txt

# ─── Étape 5 : Modules et bibliothèques chargés ───────────────────────────────
$V3 -f $DUMP linux.lsmod > analysis/lsmod.txt
# Comparer avec la liste des modules attendus sur Ubuntu 22.04

# ─── Étape 6 : Historique bash en mémoire ────────────────────────────────────
$V3 -f $DUMP linux.bash > analysis/bash_history_mem.txt
cat analysis/bash_history_mem.txt

# ─── Étape 7 : Recherche de clés et chaînes suspectes ────────────────────────
$V3 -f $DUMP linux.yarascan \
  --yara-rules '/tmp/rules.yar' > analysis/yara_matches.txt
# Règle YARA minimale (fichier /tmp/rules.yar) :
# rule SuspiciousStrings {
#   strings:
#     $a = "wget " nocase
#     $b = "base64 -d" nocase
#     $c = "curl " nocase
#     $d = "nc -e" nocase
#   condition: any of them
# }
\end{lstlisting}

\begin{boxexo}
\textbf{Travail~C1~:}
\begin{enumerate}
  \item Listez tous les processus actifs au moment du dump.
        Identifiez les 3~processus les plus suspects et justifiez.
  \item Les connexions réseau révèlent-elles une connexion C2 active~?
        Si oui~: IP, port, processus.
  \item La commande \texttt{linux.bash} révèle-t-elle des commandes
        supplémentaires par rapport à l'historique sur disque~?
        Pourquoi est-ce forensiquement important~?
  \item Qu'est-ce que le \textit{module LiME} et pourquoi son
        chargement sur une machine suspecte serait un indicateur
        de compromission~?
\end{enumerate}
\end{boxexo}


% ==============================================================================
\section{Partie~D — Construction de la timeline}
\label{lab2:partieD}
% ==============================================================================

\subsection{D1 — Timeline multi-sources}

\begin{lstlisting}[style=StyleTerminal,
  caption={Construction de la timeline forensique complète},
  label={lst:lab2-timeline}]
# ─── The Sleuth Kit : bodyfile depuis l'image ─────────────────────────────────
fls -r -m "/" \
  ~/cases/INC-2026-0412-001/evidence/srv-web-01.dd > \
  ~/cases/INC-2026-0412-001/analysis/bodyfile.txt

# Générer la timeline NTFS (ici ext4) pour le 12/04/2026
mactime -b ~/cases/INC-2026-0412-001/analysis/bodyfile.txt \
        -d 2026-04-12 \
        -z UTC > \
        ~/cases/INC-2026-0412-001/analysis/timeline_20260412.csv

# Filtrer sur la fenêtre d'incident (02h00 à 06h00)
grep "^2026-04-12 0[2345]:" \
  ~/cases/INC-2026-0412-001/analysis/timeline_20260412.csv | \
  sort > ~/cases/INC-2026-0412-001/analysis/timeline_incident.csv

wc -l ~/cases/INC-2026-0412-001/analysis/timeline_incident.csv

# ─── Croiser avec les logs auth.log ───────────────────────────────────────────
grep "Apr 12 0[2345]" \
  /mnt/forensics/srv-web-01/var/log/auth.log | \
  awk '{printf "2026-04-12 %s:%s:%s | auth.log | %s\n", $3, $4, $5, $0}' >> \
  ~/cases/INC-2026-0412-001/analysis/timeline_incident.csv

# Trier la timeline finale
sort ~/cases/INC-2026-0412-001/analysis/timeline_incident.csv > \
  ~/cases/INC-2026-0412-001/analysis/timeline_FINAL.csv
\end{lstlisting}

\begin{boxexo}
\textbf{Travail~D1~:} Construisez une timeline forensique au format
tableau (8 colonnes minimum)~: timestamp UTC, source (auth.log,
filesystem, mémoire), entité (processus/fichier/utilisateur),
type d'événement, chemin/valeur, IP associée, TTP MITRE,
niveau de confiance (Confirmé/Probable/Suspect).

La timeline doit contenir au minimum 10~entrées sur la plage 02h00--06h00.
Annotez chaque entrée avec le TTP MITRE correspondant.
\end{boxexo}


% ==============================================================================
\section{Partie~E — Rapport d'investigation forensique}
\label{lab2:partieE}
% ==============================================================================

\textbf{Rédigez un rapport d'investigation forensique complet}
en respectant la structure ci-dessous (20~pages maximum, Times New Roman
12, interligne 1,5)~:

\begin{enumerate}
  \item \textbf{Page de garde} (mandat, pièces saisies, investigateur,
        dates).
  \item \textbf{Résumé exécutif} (1~page, non technique, accessible
        au directeur).
  \item \textbf{Déclaration d'intégrité} (hashes vérifiés,
        méthodes utilisées, conformité ISO~27037).
  \item \textbf{Environnement d'analyse} (OS, outils, versions).
  \item \textbf{Chronologie de l'incident} (timeline annotée).
  \item \textbf{Artefacts découverts} (tableau complet avec hashes).
  \item \textbf{Analyse et interprétation} (reconstruction de l'attaque,
        TTPs MITRE, vecteur initial, propagation).
  \item \textbf{Réponses aux questions posées}~:
        (a)~Comment l'attaquant a-t-il accédé au serveur~?
        (b)~Quelles actions a-t-il menées~?
        (c)~Des données ont-elles été exfiltrées~?
        (d)~Des persistances ont-elles été installées~?
  \item \textbf{Recommandations} (5~mesures correctives priorisées).
  \item \textbf{Annexes} (copies de commandes, captures d'écran, logs).
\end{enumerate}

\begin{table}[H]
\centering\small
\caption{Barème Lab~2}
\label{tab:lab2-bareme}
\rowcolors{2}{TableauPair}{TableauImpair}
\begin{tabular}{p{6cm}p{7cm}p{2cm}}
\toprule
\tableauHeader{Critère} & \tableauHeader{Indicateurs} &
\tableauHeader{Points} \\
\midrule
Intégrité vérifiée et documentée & Hashes corrects, PV analysé & 1 \\
Montage correct et traces documentées & ro confirmé, commandes loguées & 1 \\
Artefacts identifiés (disque) & $\geq 5$ artefacts pertinents & 2 \\
Analyse mémoire & $\geq 3$ observations Volatility pertinentes & 2 \\
Timeline complète et annotée MITRE & $\geq 10$ entrées, TTPs corrects & 2 \\
Qualité du rapport & Structure, langage, conclusion & 2 \\
\bottomrule
\multicolumn{2}{r}{\textbf{Total}} & \textbf{10} \\
\end{tabular}
\end{table}
